%==============================================================================
%  k-Means Clustering:
%  Background, Theory, Types and Worked Examples
%
%  Self-contained lecture notes for an introductory Machine Learning course
%  (B.Tech / M.Tech level).  Prerequisites: basic linear algebra and distances.
%
%  Compile with:  pdflatex k-means_clustering.tex   (run twice for refs)
%  Requires a full TeX distribution (TeX Live / MiKTeX) with pgfplots + tikz.
%==============================================================================
\documentclass[11pt,a4paper]{report}

%------------------------------- Page geometry --------------------------------
\usepackage[a4paper,margin=1in]{geometry}

%------------------------------- Mathematics ----------------------------------
\usepackage{amsmath,amssymb,amsthm,mathtools}
\usepackage{bm}

%------------------------------- Tables & layout ------------------------------
\usepackage{booktabs}
\usepackage{array}
\usepackage{multirow}
\usepackage{enumitem}
\usepackage{caption}

%------------------------------- Graphics -------------------------------------
\usepackage{graphicx}
\usepackage{xcolor}
\usepackage{tikz}
\usepackage{pgfplots}
\pgfplotsset{compat=1.18}
\usetikzlibrary{arrows.meta,calc}

%------------------------------- Code listings --------------------------------
\usepackage{listings}

%------------------------------- Hyperlinks (load last) -----------------------
\usepackage[colorlinks=true,linkcolor=blue!60!black,citecolor=blue!60!black,
urlcolor=blue!60!black,linktocpage=true]{hyperref}

%------------------------------- Theorem-like envs ----------------------------
\theoremstyle{definition}
\newtheorem{definition}{Definition}[chapter]
\newtheorem{example}{Example}[chapter]
\newtheorem{exercise}{Exercise}[chapter]
\theoremstyle{plain}
\newtheorem{theorem}{Theorem}[chapter]
\theoremstyle{remark}
\newtheorem*{remark}{Remark}

%------------------------------- Handy macros ---------------------------------
\newcommand{\R}{\mathbb{R}}
\newcommand{\vx}{\mathbf{x}}
\newcommand{\vy}{\mathbf{y}}
\newcommand{\vmu}{\bm{\mu}}
\newcommand{\C}{\mathcal{C}}
\newcommand{\norm}[1]{\left\lVert #1 \right\rVert}
\DeclareMathOperator*{\argmin}{arg\,min}
\DeclareMathOperator*{\argmax}{arg\,max}

%------------------------------- Listing style --------------------------------
\definecolor{codegreen}{rgb}{0.0,0.5,0.0}
\definecolor{codegray}{rgb}{0.5,0.5,0.5}
\definecolor{codeblue}{rgb}{0.1,0.1,0.6}
\definecolor{backcol}{rgb}{0.97,0.97,0.95}
\lstdefinestyle{python}{
	backgroundcolor=\color{backcol},language=Python,
	commentstyle=\color{codegreen},keywordstyle=\color{codeblue}\bfseries,
	numberstyle=\tiny\color{codegray},stringstyle=\color{codegreen},
	basicstyle=\ttfamily\footnotesize,breaklines=true,captionpos=b,
	numbers=left,numbersep=6pt,showstringspaces=false,frame=single,
	rulecolor=\color{codegray},tabsize=2}

%==============================================================================
\begin{document}
	
	%------------------------------- Title ----------------------------------------
	\begin{titlepage}
		\centering
		\vspace*{2cm}
		{\Huge\bfseries $k$-Means Clustering\par}
		\vspace{0.6cm}
		{\Large Background, Theory, Types\\[2pt]
			and Worked Examples by Hand\par}
		\vspace{1.5cm}
		{\large A Self-Contained Set of Lecture Notes\par}
		\vspace{0.4cm}
		{\large Introduction to Machine Learning\par}
		\vspace{2.5cm}
		\begin{tabular}{c}
			\hline\\[-6pt]
			\textbf{Prerequisites:} Basic Linear Algebra and Distance Measures\\[4pt]
			\textbf{Level:} B.Tech / M.Tech\\[4pt]
			\hline
		\end{tabular}
		\vfill
		{\large \today\par}
	\end{titlepage}
	
	\tableofcontents
	\clearpage
	
	%==============================================================================
	\chapter{Background: Clustering in Machine Learning}
	%==============================================================================
	
	\section{Supervised vs.\ Unsupervised Learning}
	In \emph{supervised} learning every training example carries a label and we
	learn a mapping from inputs to labels. In \emph{unsupervised} learning we are
	given only the inputs
	\[
	\{\vx^{(1)},\vx^{(2)},\dots,\vx^{(m)}\},\qquad \vx^{(i)}\in\R^{n},
	\]
	with \textbf{no labels}, and must discover structure in the data.
	\emph{Clustering}---grouping similar objects together---is the most important
	unsupervised task, and \textbf{$k$-means} is by far the most widely used
	clustering algorithm.
	
	\section{What is Clustering?}
	\begin{definition}[Clustering]
		Clustering partitions a set of objects into groups (\emph{clusters}) so that
		objects within a cluster are similar (close together) and objects in different
		clusters are dissimilar (far apart).
	\end{definition}
	
	\section{Partitional vs.\ Hierarchical Clustering}
	\begin{center}
		\begin{tabular}{@{}p{6.4cm}p{6.4cm}@{}}
			\toprule
			\textbf{Partitional (flat)} & \textbf{Hierarchical}\\
			\midrule
			One partition into a fixed number $k$ of clusters. & A nested tree of clusters
			at all levels (a dendrogram).\\
			Must choose $k$ in advance. & $k$ need not be fixed in advance.\\
			Example: \textbf{$k$-means}, $k$-medoids. & Example: agglomerative, divisive.\\
			\bottomrule
		\end{tabular}
	\end{center}
	\textbf{$k$-means is the prototypical partitional method}: it divides the data
	into exactly $k$ non-overlapping clusters, each represented by its centre.
	
	\section{Where does $k$-Means Fit? Applications}
	$k$-means is unsupervised, partitional, and \emph{centroid-based} (each cluster
	is summarized by a mean vector). Typical uses:
	\begin{itemize}[nosep]
		\item \textbf{Customer segmentation} by purchase behaviour.
		\item \textbf{Image compression / colour quantization} (cluster pixel colours,
		replace each by its centroid).
		\item \textbf{Document grouping} and vector quantization.
		\item \textbf{Feature learning} (cluster-based codebooks) and data
		pre-processing.
		\item \textbf{Anomaly detection} (points far from every centroid).
	\end{itemize}
	
	%==============================================================================
	\chapter{Prerequisites: Distance and Centroids}
	%==============================================================================
	
	Each object is a \emph{point} in $n$-dimensional space, $\vx=(x_1,\dots,x_n)$,
	where $n$ is the number of features. $k$-means needs two ingredients: a way to
	measure how far two points are, and a way to summarize a group of points.
	
	\section{Euclidean Distance}
	For two points $\vx=(x_1,\dots,x_n)$ and $\vy=(y_1,\dots,y_n)$,
	\begin{equation}
		d(\vx,\vy)=\norm{\vx-\vy}=\sqrt{\sum_{k=1}^{n}(x_k-y_k)^2}.
		\label{eq:euclid}
	\end{equation}
	In the plane ($n=2$) this is the ordinary
	$\sqrt{(x_1-y_1)^2+(x_2-y_2)^2}$. Here $\vx$ and $\vy$ are two \emph{different
		points}, while $x_k$ is the $k$-th coordinate of the single point $\vx$ (not a
	separate point).
	
	\begin{remark}[Squared distance suffices for assignment]
		$k$-means assigns each point to its \emph{nearest} centroid. Because the square
		root is monotonic, $\argmin_j\norm{\vx-\vmu_j}=\argmin_j\norm{\vx-\vmu_j}^2$, so
		we may compare \emph{squared} distances and avoid square roots entirely. We do
		this in the worked examples.
	\end{remark}
	
	\section{Centroid (Mean) of a Cluster}
	\begin{definition}[Centroid]
		The centroid of a cluster $\C=\{\vx^{(1)},\dots,\vx^{(p)}\}$ is the coordinate-wise
		mean
		\begin{equation}
			\vmu=\frac{1}{p}\sum_{i=1}^{p}\vx^{(i)}
			=\Bigl(\tfrac1p\textstyle\sum_i x^{(i)}_1,\ \dots,\ \tfrac1p\sum_i x^{(i)}_n\Bigr).
			\label{eq:centroid}
		\end{equation}
	\end{definition}
	The centroid is the point that \emph{minimizes} the sum of squared distances to
	the members of $\C$ (proved in Chapter~\ref{ch:theory}); this is precisely why
	$k$-means updates each centre to the mean.
	
	\begin{remark}[Feature scaling]
		Euclidean distance is dominated by large-range features. Standardize each
		feature (zero mean, unit variance) before running $k$-means so that no single
		feature dominates the clustering.
	\end{remark}
	
	%==============================================================================
	\chapter{Theory of $k$-Means}
	\label{ch:theory}
	%==============================================================================
	
	\section{The Objective Function (WCSS / Inertia)}
	$k$-means seeks $k$ centroids $\vmu_1,\dots,\vmu_k$ and an assignment of each
	point to one cluster that \emph{minimize the within-cluster sum of squares}
	(WCSS), also called the \emph{inertia} or distortion:
	\begin{equation}
		\boxed{\;
			J(\C_1,\dots,\C_k;\vmu_1,\dots,\vmu_k)
			=\sum_{j=1}^{k}\ \sum_{\vx\in\C_j}\norm{\vx-\vmu_j}^2 .
			\;}
		\label{eq:wcss}
	\end{equation}
	$J$ is the total squared distance from every point to the centre of its cluster.
	Small $J$ means tight, compact clusters. Minimizing $J$ exactly is NP-hard, so
	we use an efficient iterative heuristic---\emph{Lloyd's algorithm}---that is
	what everyone means by ``$k$-means''.
	
	\section{Lloyd's Algorithm}
	The algorithm alternates two steps that each \emph{decrease} $J$.
	
	\begin{center}
		\fbox{\parbox{0.92\textwidth}{
				\textbf{$k$-Means (Lloyd's Algorithm)}
				\begin{enumerate}[nosep,leftmargin=2.4em]
					\item \textbf{Initialize.} Choose $k$ and pick $k$ initial centroids
					$\vmu_1,\dots,\vmu_k$ (e.g.\ $k$ random data points).
					\item \textbf{Assignment step.} Assign every point to the cluster of its
					nearest centroid:
					\[
					\C_j=\bigl\{\vx : j=\argmin_{\ell}\ \norm{\vx-\vmu_\ell}^2\bigr\}.
					\]
					\item \textbf{Update step.} Recompute each centroid as the mean of the points
					now assigned to it: $\displaystyle \vmu_j=\frac{1}{|\C_j|}\sum_{\vx\in\C_j}\vx$.
					\item \textbf{Repeat} steps 2--3 until the assignments (equivalently the
					centroids) stop changing.
				\end{enumerate}
		}}
	\end{center}
	
	\section{Why Each Step Decreases $J$ (Convergence)}
	Both steps minimize \eqref{eq:wcss} over one set of variables while holding the
	other fixed---this is \emph{coordinate descent}.
	\begin{itemize}[nosep]
		\item \textbf{Assignment step} fixes the centroids and minimizes $J$ over the
		assignments. Sending each point to its nearest centroid makes its
		contribution $\norm{\vx-\vmu_j}^2$ as small as possible, so $J$ cannot
		increase.
		\item \textbf{Update step} fixes the assignments and minimizes $J$ over the
		centroids. For a fixed cluster $\C_j$, the minimizer of
		$\sum_{\vx\in\C_j}\norm{\vx-\vmu}^2$ is the mean:
		\[
		\frac{\partial}{\partial\vmu}\sum_{\vx\in\C_j}\norm{\vx-\vmu}^2
		=-2\sum_{\vx\in\C_j}(\vx-\vmu)=\mathbf 0
		\ \Longrightarrow\ \vmu=\frac{1}{|\C_j|}\sum_{\vx\in\C_j}\vx .
		\]
	\end{itemize}
	Because $J$ never increases and there are only finitely many possible
	assignments, the algorithm \textbf{converges} in a finite number of iterations.
	
	\begin{remark}[Local, not global, optimum]
		$J$ is non-convex, so Lloyd's algorithm converges to a \emph{local} minimum that
		depends on the initialization. In practice one runs it several times with
		different seeds and keeps the clustering with the smallest $J$ (this is the
		\texttt{n\_init} parameter in libraries).
	\end{remark}
	
	\section{Complexity}
	Each iteration computes $m\times k$ distances in $\R^n$, costing
	$\mathcal{O}(mkn)$ time. With $t$ iterations the total is $\mathcal{O}(tmkn)$,
	which is \emph{linear} in the number of points---this scalability is the main
	reason for $k$-means's popularity.
	
	%==============================================================================
	\chapter{Types and Variants of $k$-Means}
	%==============================================================================
	
	\section{Common Variants}
	\begin{description}[leftmargin=0pt,itemsep=4pt]
		\item[\textbf{Lloyd's / standard $k$-means.}] The batch algorithm of
		Chapter~\ref{ch:theory}; centroids are means, distance is Euclidean.
		
		\item[\textbf{$k$-means\texttt{++}.}] A smarter \emph{initialization}: seeds are
		chosen one at a time with probability proportional to squared distance from the
		already-chosen seeds. This spreads the initial centres out, giving faster and
		more reliable convergence. It is the default initializer in most libraries.
		
		\item[\textbf{Mini-batch $k$-means.}] Updates centroids using small random
		batches of data instead of the whole dataset each iteration. Much faster on
		very large / streaming data, at a small cost in quality.
		
		\item[\textbf{$k$-medoids (PAM).}] Cluster centres are \emph{actual data points}
		(medoids) rather than means. More robust to outliers and works with any
		dissimilarity (not just Euclidean), but costlier.
		
		\item[\textbf{$k$-medians.}] Uses the median (and $L_1$/Manhattan distance)
		instead of the mean; more robust to outliers than $k$-means.
		
		\item[\textbf{Fuzzy $c$-means (soft $k$-means).}] Each point gets a
		\emph{membership degree} in every cluster (values in $[0,1]$ summing to one)
		rather than a hard single assignment.
	\end{description}
	
	\section{Choosing the Number of Clusters $k$}
	$k$ must be supplied. Two standard heuristics:
	\begin{itemize}[nosep]
		\item \textbf{Elbow method.} Plot WCSS $J$ against $k$. $J$ always decreases
		as $k$ grows, but the rate of decrease drops sharply at the ``right''
		$k$, producing an elbow. Pick $k$ at the elbow
		(worked in Chapter~\ref{ch:elbow}).
		\item \textbf{Silhouette score.} For each point compare its average distance
		to its own cluster ($a$) with the average distance to the nearest other
		cluster ($b$): the silhouette is $\tfrac{b-a}{\max(a,b)}\in[-1,1]$.
		Choose the $k$ that maximizes the mean silhouette.
	\end{itemize}
	
	%==============================================================================
	\chapter{Worked Example: $k$-Means by Hand}
	\label{ch:worked}
	%==============================================================================
	
	\section{The Data and Setup}
	We cluster six points in the plane with $k=2$.
	\begin{center}
		\begin{tabular}{@{}ccccccc@{}}
			\toprule
			Point & $P_1$ & $P_2$ & $P_3$ & $P_4$ & $P_5$ & $P_6$\\
			\midrule
			$(x,y)$ & $(1,1)$ & $(1,2)$ & $(2,2)$ & $(4,3)$ & $(5,4)$ & $(4,4)$\\
			\bottomrule
		\end{tabular}
	\end{center}
	To show the algorithm actually doing work, we deliberately pick a \emph{poor}
	initialization: both seeds inside the left group,
	\[
	\vmu_1=P_1=(1,1),\qquad \vmu_2=P_2=(1,2).
	\]
	We assign by \emph{squared} Euclidean distance
	$d^2=(x-\mu_x)^2+(y-\mu_y)^2$ (nearest centroid wins; ties go to $C_1$).
	
	\section{Iteration 1}
	\subsection*{Assignment (centroids $\vmu_1=(1,1),\ \vmu_2=(1,2)$)}
	\begin{center}
		\begin{tabular}{@{}cccc@{}}
			\toprule
			Point & $d^2$ to $\vmu_1$ & $d^2$ to $\vmu_2$ & Assign\\
			\midrule
			$P_1(1,1)$ & $0$  & $1$  & $C_1$\\
			$P_2(1,2)$ & $1$  & $0$  & $C_2$\\
			$P_3(2,2)$ & $2$  & $1$  & $C_2$\\
			$P_4(4,3)$ & $13$ & $10$ & $C_2$\\
			$P_5(5,4)$ & $25$ & $20$ & $C_2$\\
			$P_6(4,4)$ & $18$ & $13$ & $C_2$\\
			\bottomrule
		\end{tabular}
	\end{center}
	So $C_1=\{P_1\}$ and $C_2=\{P_2,P_3,P_4,P_5,P_6\}$.
	
	\subsection*{Update (new centroids = cluster means)}
	\[
	\vmu_1=(1,1),\qquad
	\vmu_2=\Bigl(\tfrac{1+2+4+5+4}{5},\tfrac{2+2+3+4+4}{5}\Bigr)
	=\Bigl(\tfrac{16}{5},\tfrac{15}{5}\Bigr)=(3.2,\,3.0).
	\]
	The WCSS after this update is
	\[
	J=\underbrace{0}_{C_1}
	+\underbrace{5.84+2.44+0.64+4.24+1.64}_{C_2\ \text{(squared dists to }(3.2,3))}
	=14.8 .
	\]
	
	\section{Iteration 2}
	\subsection*{Assignment (centroids $\vmu_1=(1,1),\ \vmu_2=(3.2,3)$)}
	\begin{center}
		\begin{tabular}{@{}cccc@{}}
			\toprule
			Point & $d^2$ to $\vmu_1=(1,1)$ & $d^2$ to $\vmu_2=(3.2,3)$ & Assign\\
			\midrule
			$P_1(1,1)$ & $0$  & $2.2^2+2^2=8.84$   & $C_1$\\
			$P_2(1,2)$ & $1$  & $2.2^2+1^2=5.84$   & $C_1$\\
			$P_3(2,2)$ & $2$  & $1.2^2+1^2=2.44$   & $C_1$\\
			$P_4(4,3)$ & $13$ & $0.8^2+0^2=0.64$   & $C_2$\\
			$P_5(5,4)$ & $25$ & $1.8^2+1^2=4.24$   & $C_2$\\
			$P_6(4,4)$ & $18$ & $0.8^2+1^2=1.64$   & $C_2$\\
			\bottomrule
		\end{tabular}
	\end{center}
	Now $P_2$ and $P_3$ have \textbf{switched} to $C_1$:
	$C_1=\{P_1,P_2,P_3\}$, $C_2=\{P_4,P_5,P_6\}$.
	
	\subsection*{Update}
	\[
	\vmu_1=\Bigl(\tfrac{1+1+2}{3},\tfrac{1+2+2}{3}\Bigr)=\Bigl(\tfrac43,\tfrac53\Bigr)
	\approx(1.333,1.667),\qquad
	\vmu_2=\Bigl(\tfrac{4+5+4}{3},\tfrac{3+4+4}{3}\Bigr)=\Bigl(\tfrac{13}{3},\tfrac{11}{3}\Bigr)
	\approx(4.333,3.667).
	\]
	The new WCSS is
	\[
	J=\tfrac43+\tfrac43=\tfrac83\approx 2.667,
	\]
	a large drop from $14.8$ (each cluster contributes $\tfrac43$; e.g.\ for $C_1$:
	$\tfrac59+\tfrac29+\tfrac59=\tfrac43$).
	
	\section{Iteration 3 --- Convergence Check}
	Re-running the assignment with $\vmu_1\approx(1.333,1.667)$ and
	$\vmu_2\approx(4.333,3.667)$ leaves every point with its \emph{nearest} centroid
	unchanged ($P_1,P_2,P_3\!\to\!C_1$; $P_4,P_5,P_6\!\to\!C_2$). Since no point
	changes cluster, the centroids do not move: the algorithm has
	\textbf{converged}. Final result:
	\[
	C_1=\{P_1,P_2,P_3\},\quad \vmu_1=(\tfrac43,\tfrac53);\qquad
	C_2=\{P_4,P_5,P_6\},\quad \vmu_2=(\tfrac{13}{3},\tfrac{11}{3});\qquad J^\star=\tfrac83 .
	\]
	
	\begin{figure}[ht]
		\centering
		\begin{tikzpicture}
			\begin{axis}[
				width=11cm,height=7.5cm,
				xlabel={$x$},ylabel={$y$},
				xmin=0,xmax=6,ymin=0,ymax=5,grid=both,
				legend pos=north west,legend cell align=left]
				% cluster 1 points
				\addplot[only marks,mark=*,mark size=2.6pt,color=blue]
				coordinates {(1,1)(1,2)(2,2)};
				\addlegendentry{$C_1=\{P_1,P_2,P_3\}$}
				% cluster 2 points
				\addplot[only marks,mark=square*,mark size=2.6pt,color=red!80!black]
				coordinates {(4,3)(5,4)(4,4)};
				\addlegendentry{$C_2=\{P_4,P_5,P_6\}$}
				% centroids
				\addplot[only marks,mark=star,mark size=4.5pt,color=black]
				coordinates {(1.3333,1.6667)(4.3333,3.6667)};
				\addlegendentry{centroids $\bm\mu_1,\bm\mu_2$}
			\end{axis}
		\end{tikzpicture}
		\caption{Converged $k$-means clustering ($k=2$). The two centroids (stars) sit
			at the mean of their assigned points; WCSS $=8/3$.}
		\label{fig:kmeans}
	\end{figure}
	
	\section{Summary of the Run}
	\begin{center}
		\begin{tabular}{@{}cccc@{}}
			\toprule
			Iteration & $C_1$ & $C_2$ & WCSS $J$\\
			\midrule
			1 & $\{P_1\}$ & $\{P_2,P_3,P_4,P_5,P_6\}$ & $14.800$\\
			2 & $\{P_1,P_2,P_3\}$ & $\{P_4,P_5,P_6\}$ & $2.667$\\
			3 & $\{P_1,P_2,P_3\}$ & $\{P_4,P_5,P_6\}$ & $2.667$ (converged)\\
			\bottomrule
		\end{tabular}
	\end{center}
	Even from a deliberately bad start, $k$-means recovered the natural two groups
	in a couple of iterations, with $J$ decreasing monotonically
	$14.8\to2.667\to2.667$, exactly as the theory guarantees.
	
	%==============================================================================
	\chapter{Worked Example: A Larger, Slower Run (5 Iterations)}
	\label{ch:worked2}
	%==============================================================================
	
	The previous example converged in just two useful iterations. To see the
	centroids \emph{crawl} gradually into place, we now cluster \textbf{twelve}
	points arranged as a gently rising \emph{chain}, with $k=2$ from a deliberately
	poor start---both seeds at the lower-left end. The cluster boundary then
	advances one point per pass, and the run takes \textbf{five iterations} to
	converge.
	
	\section{The Data and Setup}
	\begin{center}
	\begin{tabular}{@{}c|cccccc@{}}
	\toprule
	Point & $P_1$ & $P_2$ & $P_3$ & $P_4$ & $P_5$ & $P_6$\\
	$(x,y)$ & $(1,1)$ & $(2,1)$ & $(3,2)$ & $(4,2)$ & $(5,3)$ & $(6,3)$\\
	\midrule
	Point & $P_7$ & $P_8$ & $P_9$ & $P_{10}$ & $P_{11}$ & $P_{12}$\\
	$(x,y)$ & $(7,4)$ & $(8,4)$ & $(9,5)$ & $(10,5)$ & $(11,6)$ & $(12,6)$\\
	\bottomrule
	\end{tabular}
	\end{center}
	Initial centroids (both at the lower-left end):
	$\vmu_1=(1,1)$, $\vmu_2=(2,1)$. As before we assign by squared distance
	$d^2=(x-\mu_x)^2+(y-\mu_y)^2$ (nearest centroid wins; ties go to $C_1$).
	
	\section{Iteration 1 in Detail}
	\begin{center}
	\begin{tabular}{@{}ccccc@{}}
	\toprule
	Point & $(x,y)$ & $d^2$ to $\vmu_1{=}(1,1)$ & $d^2$ to $\vmu_2{=}(2,1)$ & Assign\\
	\midrule
	$P_1$ & $(1,1)$ & $0$ & $1$ & $C_1$\\
	$P_2$ & $(2,1)$ & $1$ & $0$ & $C_2$\\
	$P_3$ & $(3,2)$ & $5$ & $2$ & $C_2$\\
	$P_4$ & $(4,2)$ & $10$ & $5$ & $C_2$\\
	$P_5$ & $(5,3)$ & $20$ & $13$ & $C_2$\\
	$P_6$ & $(6,3)$ & $29$ & $20$ & $C_2$\\
	$P_7$ & $(7,4)$ & $45$ & $34$ & $C_2$\\
	$P_8$ & $(8,4)$ & $58$ & $45$ & $C_2$\\
	$P_9$ & $(9,5)$ & $80$ & $65$ & $C_2$\\
	$P_{10}$ & $(10,5)$ & $97$ & $80$ & $C_2$\\
	$P_{11}$ & $(11,6)$ & $125$ & $106$ & $C_2$\\
	$P_{12}$ & $(12,6)$ & $146$ & $125$ & $C_2$\\
	\bottomrule
	\end{tabular}
	\end{center}
	So $C_1=\{P_1\}$ and $C_2=\{P_2,\dots,P_{12}\}$. Updating (the mean of the $11$
	points now in $C_2$),
	\[
	\vmu_1=(1,1),\qquad
	\vmu_2=\Bigl(\tfrac{77}{11},\tfrac{41}{11}\Bigr)=(7.0,\,3.727),
	\]
	with WCSS $J=138.18$. The right-hand centroid leaps to the middle of the chain,
	which will pull the next several points across one at a time.
	
	\section{Full Evolution (Iterations 1--5)}
	Each row shows the centroids used for the assignment, the resulting membership,
	and the updated centroids/WCSS; the updated centroids become the next row's
	starting centroids.
	\begin{center}
	\small
	\setlength{\tabcolsep}{4pt}
	\begin{tabular}{@{}c cc l l c@{}}
	\toprule
	It & $\vmu_1$ (start) & $\vmu_2$ (start) & $C_1$ & $C_2$ & $J$\\
	\midrule
	1 & $(1,1)$      & $(2,1)$        & $P_1$          & $P_2\!-\!P_{12}$ & $138.18$\\
	2 & $(1,1)$      & $(7.0,3.727)$  & $P_1\!-\!P_4$   & $P_5\!-\!P_{12}$ & $58.00$\\
	3 & $(2.5,1.5)$  & $(8.5,4.5)$    & $P_1\!-\!P_5$   & $P_6\!-\!P_{12}$ & $48.23$\\
	4 & $(3.0,1.8)$  & $(9.0,4.714)$  & $P_1\!-\!P_6$   & $P_7\!-\!P_{12}$ & $43.00$\\
	5 & $(3.5,2.0)$  & $(9.5,5.0)$    & $P_1\!-\!P_6$   & $P_7\!-\!P_{12}$ & $43.00$\\
	\bottomrule
	\end{tabular}
	\end{center}
	At iteration~5 the membership (and hence the centroids) is identical to
	iteration~4, so the algorithm has \textbf{converged}. The final clustering is
	\[
	C_1=\{P_1,\dots,P_6\},\ \vmu_1=\bigl(\tfrac{7}{2},2\bigr)=(3.5,2.0);
	\quad
	C_2=\{P_7,\dots,P_{12}\},\ \vmu_2=\bigl(\tfrac{19}{2},5\bigr)=(9.5,5.0);
	\]
	with minimum WCSS $J^\star=43.0$.
	
	\section{Discussion and Figure}
	Watch the left cluster grow one point at a time: $C_1$ starts as just $\{P_1\}$,
	then gains $P_2,P_3,P_4$ at iteration~2, $P_5$ at iteration~3, and $P_6$ at
	iteration~4, after which nothing changes. Each new member drags $\vmu_1$ up the
	chain while $\vmu_2$ retreats toward the top-right points. The cost falls
	monotonically
	\[
	138.18\;\to\;58.00\;\to\;48.23\;\to\;43.00\;\to\;43.00,
	\]
	never increasing---the convergence guarantee in action. This gradual,
	five-iteration crawl is typical when the data form a continuum and the initial
	centroids are placed badly; a spread-out start (e.g.\ $k$-means\texttt{++})
	would reach the same split in fewer steps.
	
	\begin{figure}[ht]
	\centering
	\begin{tikzpicture}
	\begin{axis}[
	  width=12cm,height=7cm,
	  xlabel={$x$},ylabel={$y$},
	  xmin=0,xmax=13,ymin=0,ymax=7,grid=both,
	  legend pos=north west,legend cell align=left]
	% final cluster 1
	\addplot[only marks,mark=*,mark size=2.4pt,color=blue]
	  coordinates {(1,1)(2,1)(3,2)(4,2)(5,3)(6,3)};
	\addlegendentry{$C_1=\{P_1,\dots,P_6\}$}
	% final cluster 2
	\addplot[only marks,mark=square*,mark size=2.4pt,color=red!80!black]
	  coordinates {(7,4)(8,4)(9,5)(10,5)(11,6)(12,6)};
	\addlegendentry{$C_2=\{P_7,\dots,P_{12}\}$}
	% centroid 1 trajectory
	\addplot[thick,black,mark=o,mark size=2pt,dashed]
	  coordinates {(1,1)(2.5,1.5)(3,1.8)(3.5,2)};
	\addlegendentry{$\bm\mu_1$ path}
	% centroid 2 trajectory
	\addplot[thick,black!55,mark=triangle,mark size=2.4pt,dashed]
	  coordinates {(2,1)(7,3.727)(8.5,4.5)(9,4.714)(9.5,5)};
	\addlegendentry{$\bm\mu_2$ path}
	% final centroids
	\addplot[only marks,mark=star,mark size=4.5pt,color=black]
	  coordinates {(3.5,2)(9.5,5)};
	\end{axis}
	\end{tikzpicture}
	\caption{The larger run over five iterations on a chain of twelve points. Dashed
	curves trace how the two centroids \emph{crawl} from their poor lower-left start
	to their final positions (stars); the final clusters are shown in blue and red.}
	\label{fig:kmeans2}
	\end{figure}
	
	%==============================================================================
	\chapter{Worked Example: Three Clusters ($k=3$) by Hand}
	\label{ch:worked3}
	%==============================================================================
	
	This example clusters \textbf{eighteen} points into $k=3$ groups, starting from
	three \emph{distinct} centroids. It also highlights a common pitfall.
	
	\begin{remark}[Initial centroids must be distinct]
	If the three seeds coincide (e.g.\ $\vmu_1=\vmu_2=\vmu_3=(0,0)$), then every
	point is \emph{equidistant} to all three centroids---a three-way tie for all
	$18$ points. Standard $k$-means then places every point in one cluster and
	leaves the other two empty, and it stalls. The seeds must therefore be
	\textbf{distinct}, as below.
	\end{remark}
	
	\section{The Data and Setup}
	\begin{center}
	\begin{tabular}{@{}c|cccccc@{}}
	\toprule
	Point & $P_1$ & $P_2$ & $P_3$ & $P_4$ & $P_5$ & $P_6$\\
	$(x,y)$ & $(3,-1)$ & $(3,0)$ & $(3,1)$ & $(4,2)$ & $(5,0)$ & $(5,1)$\\
	\midrule
	Point & $P_7$ & $P_8$ & $P_9$ & $P_{10}$ & $P_{11}$ & $P_{12}$\\
	$(x,y)$ & $(-1,3)$ & $(-1,4)$ & $(-2,2)$ & $(-2,5)$ & $(-3,2)$ & $(-3,3)$\\
	\midrule
	Point & $P_{13}$ & $P_{14}$ & $P_{15}$ & $P_{16}$ & $P_{17}$ & $P_{18}$\\
	$(x,y)$ & $(-2,-3)$ & $(-3,-3)$ & $(-3,-5)$ & $(-4,-2)$ & $(-5,-3)$ & $(-5,-4)$\\
	\bottomrule
	\end{tabular}
	\end{center}
	Initial centroids: $\vmu_1=(0,0)$, $\vmu_2=(1,1)$, $\vmu_3=(1,-1)$. We assign by
	squared distance $d^2$; on a tie the lowest cluster index wins.
	
	\section{Iteration 1 in Detail}
	The bold entry in each row is the nearest (smallest) squared distance.
	\begin{center}
	\small
	\begin{tabular}{@{}ccccc c@{}}
	\toprule
	Point & $(x,y)$ & $d^2\!\to\vmu_1(0,0)$ & $d^2\!\to\vmu_2(1,1)$ & $d^2\!\to\vmu_3(1,{-}1)$ & Assign\\
	\midrule
	$P_1$ & $(3,-1)$ & $10$ & $8$ & $\mathbf{4}$ & $C_3$\\
	$P_2$ & $(3,0)$ & $9$ & $\mathbf{5}$ & $5$ & $C_2$\\
	$P_3$ & $(3,1)$ & $10$ & $\mathbf{4}$ & $8$ & $C_2$\\
	$P_4$ & $(4,2)$ & $20$ & $\mathbf{10}$ & $18$ & $C_2$\\
	$P_5$ & $(5,0)$ & $25$ & $\mathbf{17}$ & $17$ & $C_2$\\
	$P_6$ & $(5,1)$ & $26$ & $\mathbf{16}$ & $20$ & $C_2$\\
	$P_7$ & $(-1,3)$ & $10$ & $\mathbf{8}$ & $20$ & $C_2$\\
	$P_8$ & $(-1,4)$ & $17$ & $\mathbf{13}$ & $29$ & $C_2$\\
	$P_9$ & $(-2,2)$ & $\mathbf{8}$ & $10$ & $18$ & $C_1$\\
	$P_{10}$ & $(-2,5)$ & $29$ & $\mathbf{25}$ & $45$ & $C_2$\\
	$P_{11}$ & $(-3,2)$ & $\mathbf{13}$ & $17$ & $25$ & $C_1$\\
	$P_{12}$ & $(-3,3)$ & $\mathbf{18}$ & $20$ & $32$ & $C_1$\\
	$P_{13}$ & $(-2,-3)$ & $\mathbf{13}$ & $25$ & $13$ & $C_1$\\
	$P_{14}$ & $(-3,-3)$ & $\mathbf{18}$ & $32$ & $20$ & $C_1$\\
	$P_{15}$ & $(-3,-5)$ & $34$ & $52$ & $\mathbf{32}$ & $C_3$\\
	$P_{16}$ & $(-4,-2)$ & $\mathbf{20}$ & $34$ & $26$ & $C_1$\\
	$P_{17}$ & $(-5,-3)$ & $\mathbf{34}$ & $52$ & $40$ & $C_1$\\
	$P_{18}$ & $(-5,-4)$ & $\mathbf{41}$ & $61$ & $45$ & $C_1$\\
	\bottomrule
	\end{tabular}
	\end{center}
	(Ties $P_2,P_5$ between $C_2,C_3$ and $P_{13}$ between $C_1,C_3$ go to the lower
	index.) The clusters and updated centroids are
	\[
	\begin{aligned}
	C_1&=\{9,11,12,13,14,16,17,18\}\Rightarrow \vmu_1=\bigl(-\tfrac{27}{8},-1\bigr)=(-3.375,-1.0),\\
	C_2&=\{2,3,4,5,6,7,8,10\}\Rightarrow \vmu_2=(2,2),\\
	C_3&=\{1,15\}\Rightarrow \vmu_3=(0,-3),
	\end{aligned}
	\]
	where indices denote point subscripts (so ``$9$'' means $P_9$).
	
	\section{Iteration 2 in Detail}
	Centroids $\vmu_1=(-3.375,-1)$, $\vmu_2=(2,2)$, $\vmu_3=(0,-3)$ (squared distances rounded to two decimals). The bold entry is the nearest centroid.
	\begin{center}
	\small
	\begin{tabular}{@{}ccccc c@{}}
	\toprule
	Point & $(x,y)$ & $d^2\!\to\vmu_1$ & $d^2\!\to\vmu_2$ & $d^2\!\to\vmu_3$ & Assign\\
	\midrule
	$P_{1}$ & $(3,-1)$ & $40.64$ & $\mathbf{10.00}$ & $13.00$ & $C_2$\\
	$P_{2}$ & $(3,0)$ & $41.64$ & $\mathbf{5.00}$ & $18.00$ & $C_2$\\
	$P_{3}$ & $(3,1)$ & $44.64$ & $\mathbf{2.00}$ & $25.00$ & $C_2$\\
	$P_{4}$ & $(4,2)$ & $63.39$ & $\mathbf{4.00}$ & $41.00$ & $C_2$\\
	$P_{5}$ & $(5,0)$ & $71.14$ & $\mathbf{13.00}$ & $34.00$ & $C_2$\\
	$P_{6}$ & $(5,1)$ & $74.14$ & $\mathbf{10.00}$ & $41.00$ & $C_2$\\
	$P_{7}$ & $(-1,3)$ & $21.64$ & $\mathbf{10.00}$ & $37.00$ & $C_2$\\
	$P_{8}$ & $(-1,4)$ & $30.64$ & $\mathbf{13.00}$ & $50.00$ & $C_2$\\
	$P_{9}$ & $(-2,2)$ & $\mathbf{10.89}$ & $16.00$ & $29.00$ & $C_1$\\
	$P_{10}$ & $(-2,5)$ & $37.89$ & $\mathbf{25.00}$ & $68.00$ & $C_2$\\
	$P_{11}$ & $(-3,2)$ & $\mathbf{9.14}$ & $25.00$ & $34.00$ & $C_1$\\
	$P_{12}$ & $(-3,3)$ & $\mathbf{16.14}$ & $26.00$ & $45.00$ & $C_1$\\
	$P_{13}$ & $(-2,-3)$ & $5.89$ & $41.00$ & $\mathbf{4.00}$ & $C_3$\\
	$P_{14}$ & $(-3,-3)$ & $\mathbf{4.14}$ & $50.00$ & $9.00$ & $C_1$\\
	$P_{15}$ & $(-3,-5)$ & $16.14$ & $74.00$ & $\mathbf{13.00}$ & $C_3$\\
	$P_{16}$ & $(-4,-2)$ & $\mathbf{1.39}$ & $52.00$ & $17.00$ & $C_1$\\
	$P_{17}$ & $(-5,-3)$ & $\mathbf{6.64}$ & $74.00$ & $25.00$ & $C_1$\\
	$P_{18}$ & $(-5,-4)$ & $\mathbf{11.64}$ & $85.00$ & $26.00$ & $C_1$\\
	\bottomrule
	\end{tabular}
	\end{center}
	Updated centroids:
	\[
C_1=\{9,11,12,14,16,17,18\}\Rightarrow\vmu_1=\bigl(-\tfrac{25}{7},-\tfrac{5}{7}\bigr)=(-3.571,-0.714),
\]
\[
C_2=\{1\text{--}8,10\}\Rightarrow\vmu_2=\bigl(\tfrac{19}{9},\tfrac{5}{3}\bigr)=(2.111,1.667),\qquad
C_3=\{13,15\}\Rightarrow\vmu_3=\bigl(-\tfrac{5}{2},-4\bigr)=(-2.5,-4).
\]
	
	\section{Iteration 3 in Detail}
	Centroids $\vmu_1=(-3.571,-0.714)$, $\vmu_2=(2.111,1.667)$, $\vmu_3=(-2.5,-4)$. The bold entry is the nearest centroid.
	\begin{center}
	\small
	\begin{tabular}{@{}ccccc c@{}}
	\toprule
	Point & $(x,y)$ & $d^2\!\to\vmu_1$ & $d^2\!\to\vmu_2$ & $d^2\!\to\vmu_3$ & Assign\\
	\midrule
	$P_{1}$ & $(3,-1)$ & $43.27$ & $\mathbf{7.90}$ & $39.25$ & $C_2$\\
	$P_{2}$ & $(3,0)$ & $43.69$ & $\mathbf{3.57}$ & $46.25$ & $C_2$\\
	$P_{3}$ & $(3,1)$ & $46.12$ & $\mathbf{1.23}$ & $55.25$ & $C_2$\\
	$P_{4}$ & $(4,2)$ & $64.69$ & $\mathbf{3.68}$ & $78.25$ & $C_2$\\
	$P_{5}$ & $(5,0)$ & $73.98$ & $\mathbf{11.12}$ & $72.25$ & $C_2$\\
	$P_{6}$ & $(5,1)$ & $76.41$ & $\mathbf{8.79}$ & $81.25$ & $C_2$\\
	$P_{7}$ & $(-1,3)$ & $20.41$ & $\mathbf{11.46}$ & $51.25$ & $C_2$\\
	$P_{8}$ & $(-1,4)$ & $28.84$ & $\mathbf{15.12}$ & $66.25$ & $C_2$\\
	$P_{9}$ & $(-2,2)$ & $\mathbf{9.84}$ & $17.01$ & $36.25$ & $C_1$\\
	$P_{10}$ & $(-2,5)$ & $35.12$ & $\mathbf{28.01}$ & $81.25$ & $C_2$\\
	$P_{11}$ & $(-3,2)$ & $\mathbf{7.69}$ & $26.23$ & $36.25$ & $C_1$\\
	$P_{12}$ & $(-3,3)$ & $\mathbf{14.12}$ & $27.90$ & $49.25$ & $C_1$\\
	$P_{13}$ & $(-2,-3)$ & $7.69$ & $38.68$ & $\mathbf{1.25}$ & $C_3$\\
	$P_{14}$ & $(-3,-3)$ & $5.55$ & $47.90$ & $\mathbf{1.25}$ & $C_3$\\
	$P_{15}$ & $(-3,-5)$ & $18.69$ & $70.57$ & $\mathbf{1.25}$ & $C_3$\\
	$P_{16}$ & $(-4,-2)$ & $\mathbf{1.84}$ & $50.79$ & $6.25$ & $C_1$\\
	$P_{17}$ & $(-5,-3)$ & $7.27$ & $72.35$ & $\mathbf{7.25}$ & $C_3$\\
	$P_{18}$ & $(-5,-4)$ & $12.84$ & $82.68$ & $\mathbf{6.25}$ & $C_3$\\
	\bottomrule
	\end{tabular}
	\end{center}
	Updated centroids:
	\[
C_1=\{9,11,12,16\}\Rightarrow\vmu_1=\bigl(-3,\tfrac{5}{4}\bigr)=(-3,1.25),\qquad
C_2=\{1\text{--}8,10\}\Rightarrow\vmu_2=\bigl(\tfrac{19}{9},\tfrac{5}{3}\bigr)=(2.111,1.667),
\]
\[
C_3=\{13,14,15,17,18\}\Rightarrow\vmu_3=\bigl(-\tfrac{18}{5},-\tfrac{18}{5}\bigr)=(-3.6,-3.6).
\]
	
	\section{Iteration 4 in Detail}
	Centroids $\vmu_1=(-3,1.25)$, $\vmu_2=(2.111,1.667)$, $\vmu_3=(-3.6,-3.6)$. The bold entry is the nearest centroid.
	\begin{center}
	\small
	\begin{tabular}{@{}ccccc c@{}}
	\toprule
	Point & $(x,y)$ & $d^2\!\to\vmu_1$ & $d^2\!\to\vmu_2$ & $d^2\!\to\vmu_3$ & Assign\\
	\midrule
	$P_{1}$ & $(3,-1)$ & $41.06$ & $\mathbf{7.90}$ & $50.32$ & $C_2$\\
	$P_{2}$ & $(3,0)$ & $37.56$ & $\mathbf{3.57}$ & $56.52$ & $C_2$\\
	$P_{3}$ & $(3,1)$ & $36.06$ & $\mathbf{1.23}$ & $64.72$ & $C_2$\\
	$P_{4}$ & $(4,2)$ & $49.56$ & $\mathbf{3.68}$ & $89.12$ & $C_2$\\
	$P_{5}$ & $(5,0)$ & $65.56$ & $\mathbf{11.12}$ & $86.92$ & $C_2$\\
	$P_{6}$ & $(5,1)$ & $64.06$ & $\mathbf{8.79}$ & $95.12$ & $C_2$\\
	$P_{7}$ & $(-1,3)$ & $\mathbf{7.06}$ & $11.46$ & $50.32$ & $C_1$\\
	$P_{8}$ & $(-1,4)$ & $\mathbf{11.56}$ & $15.12$ & $64.52$ & $C_1$\\
	$P_{9}$ & $(-2,2)$ & $\mathbf{1.56}$ & $17.01$ & $33.92$ & $C_1$\\
	$P_{10}$ & $(-2,5)$ & $\mathbf{15.06}$ & $28.01$ & $76.52$ & $C_1$\\
	$P_{11}$ & $(-3,2)$ & $\mathbf{0.56}$ & $26.23$ & $31.72$ & $C_1$\\
	$P_{12}$ & $(-3,3)$ & $\mathbf{3.06}$ & $27.90$ & $43.92$ & $C_1$\\
	$P_{13}$ & $(-2,-3)$ & $19.06$ & $38.68$ & $\mathbf{2.92}$ & $C_3$\\
	$P_{14}$ & $(-3,-3)$ & $18.06$ & $47.90$ & $\mathbf{0.72}$ & $C_3$\\
	$P_{15}$ & $(-3,-5)$ & $39.06$ & $70.57$ & $\mathbf{2.32}$ & $C_3$\\
	$P_{16}$ & $(-4,-2)$ & $11.56$ & $50.79$ & $\mathbf{2.72}$ & $C_3$\\
	$P_{17}$ & $(-5,-3)$ & $22.06$ & $72.35$ & $\mathbf{2.32}$ & $C_3$\\
	$P_{18}$ & $(-5,-4)$ & $31.56$ & $82.68$ & $\mathbf{2.12}$ & $C_3$\\
	\bottomrule
	\end{tabular}
	\end{center}
	Updated centroids: All eighteen points now sit with their nearest natural group; the next pass confirms it.
	\[
C_1=\{7\text{--}12\}\Rightarrow\vmu_1=\bigl(-2,\tfrac{19}{6}\bigr)=(-2.0,3.167),\qquad
C_2=\{1\text{--}6\}\Rightarrow\vmu_2=\bigl(\tfrac{23}{6},\tfrac{1}{2}\bigr)=(3.833,0.5),
\]
\[
C_3=\{13\text{--}18\}\Rightarrow\vmu_3=\bigl(-\tfrac{11}{3},-\tfrac{10}{3}\bigr)=(-3.667,-3.333).
\]
	
	\section{Iteration 5 --- Convergence}
	Reassigning every point to its nearest of the iteration-4 centroids leaves all
	memberships unchanged, so the centroids do not move: the algorithm has
	\textbf{converged} after five iterations.
	
	\section{Evolution of the Clusters}
	The table tracks each cluster's membership (point indices) per iteration; watch
	the groups sort themselves out from the poor near-origin start.
	\begin{center}
	\small
	\begin{tabular}{@{}clll@{}}
	\toprule
	It & $C_1$ & $C_2$ & $C_3$\\
	\midrule
	1 & $9,11,12,13,14,16,17,18$ & $2,3,4,5,6,7,8,10$ & $1,15$\\
	2 & $9,11,12,14,16,17,18$ & $1,2,3,4,5,6,7,8,10$ & $13,15$\\
	3 & $9,11,12,16$ & $1,2,3,4,5,6,7,8,10$ & $13,14,15,17,18$\\
	4 & $7,8,9,10,11,12$ & $1,2,3,4,5,6$ & $13,14,15,16,17,18$\\
	5 & \multicolumn{3}{l}{unchanged --- \textbf{converged}}\\
	\bottomrule
	\end{tabular}
	\end{center}
	
	\section{Final Result and Figure}
	\begin{center}
	\begin{tabular}{@{}lll@{}}
	\toprule
	Cluster & Members & Centroid\\
	\midrule
	$C_1$ (upper-left) & $P_7,P_8,P_9,P_{10},P_{11},P_{12}$ & $(-2,\tfrac{19}{6})\approx(-2.00,3.17)$\\
	$C_2$ (right)      & $P_1,P_2,P_3,P_4,P_5,P_6$          & $(\tfrac{23}{6},\tfrac12)\approx(3.83,0.50)$\\
	$C_3$ (lower-left) & $P_{13},P_{14},P_{15},P_{16},P_{17},P_{18}$ & $(-\tfrac{11}{3},-\tfrac{10}{3})\approx(-3.67,-3.33)$\\
	\bottomrule
	\end{tabular}
	\end{center}
	
	\begin{figure}[ht]
	\centering
	\begin{tikzpicture}
	\begin{axis}[
	  width=11.5cm,height=11cm,
	  axis lines=middle,
	  xlabel={$x$},ylabel={$y$},
	  xmin=-6.5,xmax=6.5,ymin=-6.5,ymax=6.5,
	  grid=both,legend pos=outer north east,legend cell align=left]
	\addplot[only marks,mark=square*,mark size=2.6pt,color=blue]
	  coordinates {(-1,3)(-1,4)(-2,2)(-2,5)(-3,2)(-3,3)};
	\addlegendentry{$C_1$ (upper-left)}
	\addplot[only marks,mark=*,mark size=2.6pt,color=red!80!black]
	  coordinates {(3,-1)(3,0)(3,1)(4,2)(5,0)(5,1)};
	\addlegendentry{$C_2$ (right)}
	\addplot[only marks,mark=triangle*,mark size=3pt,color=green!55!black]
	  coordinates {(-2,-3)(-3,-3)(-3,-5)(-4,-2)(-5,-3)(-5,-4)};
	\addlegendentry{$C_3$ (lower-left)}
	\addplot[only marks,mark=star,mark size=4pt,color=black]
	  coordinates {(-2,3.1667)(3.8333,0.5)(-3.6667,-3.3333)};
	\addlegendentry{final centroids}
	\addplot[only marks,mark=x,mark size=4pt,color=gray]
	  coordinates {(0,0)(1,1)(1,-1)};
	\addlegendentry{initial seeds}
	\end{axis}
	\end{tikzpicture}
	\caption{The $k=3$ result. The three initial seeds (gray $\times$) start bunched
	near the origin; after five iterations $k$-means recovers the three
	well-separated groups, with centroids marked by stars.}
	\label{fig:kmeans3}
	\end{figure}
	
	
	%==============================================================================
	\chapter{Worked Example: The Elbow Method}
	\label{ch:elbow}
	%==============================================================================
	
	Using the same six points, we run $k$-means for
	$k=1,2,3,4,5$ and record the (best) WCSS:
	\begin{center}
		\begin{tabular}{@{}cccccc@{}}
			\toprule
			$k$ & $1$ & $2$ & $3$ & $4$ & $5$\\
			\midrule
			WCSS $J$ & $22.17$ & $2.67$ & $1.83$ & $1.00$ & $0.50$\\
			\bottomrule
		\end{tabular}
	\end{center}
	Going from $k=1$ to $k=2$ slashes $J$ from $22.17$ to $2.67$; every further
	increase buys only a small improvement. The curve therefore bends sharply at
	$k=2$---the \textbf{elbow}---which is the recommended number of clusters, in
	agreement with the two visible groups.
	
	\begin{figure}[ht]
		\centering
		\begin{tikzpicture}
			\begin{axis}[
				width=11cm,height=6.5cm,
				xlabel={number of clusters $k$},ylabel={WCSS $J$},
				xmin=0.7,xmax=5.3,ymin=0,ymax=24,
				xtick={1,2,3,4,5},grid=both,mark size=2.6pt]
				\addplot[thick,blue,mark=*]
				coordinates {(1,22.17)(2,2.67)(3,1.83)(4,1.00)(5,0.50)};
				\node[anchor=west,font=\footnotesize] (el) at (axis cs:2.5,9) {elbow at $k=2$};
				\draw[->,thick,gray] (el.west) -- (axis cs:2.05,3.2);
			\end{axis}
		\end{tikzpicture}
		\caption{Elbow plot. WCSS drops steeply until $k=2$, then flattens; the elbow
			identifies $k=2$ as the appropriate number of clusters.}
		\label{fig:elbow}
	\end{figure}
	
	%==============================================================================
	\chapter{Pros, Cons, Comparison and Code}
	%==============================================================================
	
	\section{Advantages and Disadvantages}
	\begin{center}
		\begin{tabular}{@{}p{6.4cm}p{6.4cm}@{}}
			\toprule
			\textbf{Advantages} & \textbf{Disadvantages}\\
			\midrule
			Simple and fast; scales linearly, $\mathcal{O}(tmkn)$. & Must choose $k$ in
			advance.\\
			Easy to implement and interpret (centroids). & Converges only to a \emph{local}
			optimum; sensitive to initialization.\\
			Works well when clusters are compact and roughly spherical. & Assumes spherical,
			similar-size clusters; struggles with elongated or varied-density shapes.\\
			Guaranteed to converge. & Sensitive to outliers and to feature scaling.\\
			\bottomrule
		\end{tabular}
	\end{center}
	
	\section{Comparison with Hierarchical Clustering}
	\begin{center}
		\begin{tabular}{@{}p{6.4cm}p{6.4cm}@{}}
			\toprule
			\textbf{$k$-Means} & \textbf{Hierarchical}\\
			\midrule
			$k$ fixed \emph{before} running. & $k$ chosen \emph{after} building the tree.\\
			Depends on random initial centroids. & Deterministic.\\
			$\mathcal{O}(tmkn)$; scales to large data. &
			$\mathcal{O}(m^2)$--$\mathcal{O}(m^3)$; small/medium data.\\
			Output: $k$ flat clusters + centroids. & Output: a nested dendrogram.\\
			\bottomrule
		\end{tabular}
	\end{center}
	
	\section{Reproducing the Example in Python}
	\begin{lstlisting}[style=python,caption={$k$-means with scikit-learn.}]
		import numpy as np
		from sklearn.cluster import KMeans
		
		X = np.array([[1,1],[1,2],[2,2],[4,3],[5,4],[4,4]], dtype=float)
		
		km = KMeans(n_clusters=2, init='k-means++', n_init=10, random_state=0).fit(X)
		
		print("labels    :", km.labels_)            # e.g. [0 0 0 1 1 1]
		print("centroids :", km.cluster_centers_)   # ~[[1.333,1.667],[4.333,3.667]]
		print("WCSS (inertia):", km.inertia_)       # ~2.667
		
		# Elbow method: WCSS vs k
		for k in range(1, 6):
		j = KMeans(n_clusters=k, n_init=10, random_state=0).fit(X).inertia_
		print(f"k={k}: WCSS={j:.2f}")
	\end{lstlisting}
	The library reproduces the by-hand result: clusters
	$\{P_1,P_2,P_3\}$, $\{P_4,P_5,P_6\}$, centroids
	$(\tfrac43,\tfrac53)$ and $(\tfrac{13}{3},\tfrac{11}{3})$, and inertia $8/3$.
	
	%==============================================================================
	\chapter{Exercises}
	%==============================================================================
	
	\section{Solved Problems}
	
	\begin{example}[Nearest centroid]
		Assign $P=(3,3)$ to one of the centroids $\vmu_1=(1,1)$, $\vmu_2=(4,4)$.\\
		\textbf{Solution.} $d^2(P,\vmu_1)=2^2+2^2=8$, $d^2(P,\vmu_2)=1^2+1^2=2$. Since
		$2<8$, assign $P$ to cluster $2$.
	\end{example}
	
	\begin{example}[Centroid update]
		A cluster contains $(2,0),(4,2),(6,4)$. Find its new centroid.\\
		\textbf{Solution.}
		$\vmu=\bigl(\tfrac{2+4+6}{3},\tfrac{0+2+4}{3}\bigr)=(4,2).$
	\end{example}
	
	\begin{example}[WCSS of a cluster]
		Compute the within-cluster sum of squares of $\{(1,1),(1,2),(2,2)\}$ about its
		centroid.\\
		\textbf{Solution.} Centroid $=(\tfrac43,\tfrac53)$. Squared distances:
		$\tfrac59,\tfrac29,\tfrac59$; sum $=\tfrac{12}{9}=\tfrac43$.
	\end{example}
	
	\begin{example}[Effect of initialization]
		Why might two runs of $k$-means on the same data give different clusters?\\
		\textbf{Solution.} $J$ is non-convex; Lloyd's algorithm converges to a local
		minimum that depends on the random initial centroids. Different seeds can lead
		to different local minima---hence \texttt{n\_init} with the best-$J$ kept.
	\end{example}
	
	\section{Unsolved Problems}
	
	\begin{exercise}[One full iteration]
		Points $(1,1),(2,1),(4,3),(5,4)$, $k=2$, initial centroids $(1,1)$ and $(5,4)$.
		Perform one assignment and one update step; report the new centroids.
	\end{exercise}
	
	\begin{exercise}[Run to convergence]
		Continue Exercise~8.1 until convergence and compute the final WCSS.
	\end{exercise}
	
	\begin{exercise}[Elbow]
		For the dataset of Chapter~\ref{ch:worked}, verify by hand the WCSS values for
		$k=1$ (all points in one cluster) and $k=3$, and state which $k$ the elbow rule
		selects.
	\end{exercise}
	
	\begin{exercise}[Proof]
		Show that the mean minimizes $\sum_{\vx\in\C}\norm{\vx-\vmu}^2$ by setting the
		gradient with respect to $\vmu$ to zero.
	\end{exercise}
	
	\begin{exercise}[Squared vs Euclidean]
		Prove that assigning a point to the nearest centroid gives the same result
		whether one compares Euclidean distances or squared Euclidean distances.
	\end{exercise}
	
	\begin{exercise}[Conceptual]
		Give one dataset shape on which $k$-means fails badly (e.g.\ two concentric
		rings) and explain why, referring to the spherical-cluster assumption.
	\end{exercise}
	
	\begin{exercise}[$k$-medoids]
		State one advantage of $k$-medoids over $k$-means when the data contains
		outliers, and explain why using an actual data point as the centre helps.
	\end{exercise}
	
	%==============================================================================
	\appendix
	\chapter{Notation Summary}
	%==============================================================================
	\begin{center}
		\begin{tabular}{@{}ll@{}}
			\toprule
			Symbol & Meaning\\
			\midrule
			$m$ & number of data points\\
			$n$ & number of features (dimensions)\\
			$k$ & number of clusters (chosen in advance)\\
			$\vx^{(i)}\in\R^n$ & the $i$-th data point\\
			$\C_j$ & the $j$-th cluster; $|\C_j|$ its size\\
			$\vmu_j$ & centroid (mean) of cluster $\C_j$\\
			$d(\vx,\vy)=\norm{\vx-\vy}$ & Euclidean distance between two points\\
			$J$ & within-cluster sum of squares (WCSS / inertia)\\
			\bottomrule
		\end{tabular}
	\end{center}
	
\end{document}
